\section{Alignment as a relational operation}
\label{sec:postgres-custom-scan-scope}

The alignment pipeline described in the preceding chapters is not naturally
expressed as a scalar SQL function. Its efficient unit of work is a batch of
sequence pairs, and one pair may produce several accepted alignments. A scalar
predicate instead receives one row, returns one Boolean value, and is normally
evaluated independently at each call site. It cannot retain a group of inputs
for vector or accelerator execution, and it cannot expand one input row into a
set of match rows.

BioSeqDB nevertheless presents search through ordinary typed SQL. A typical
query is

\begin{lstlisting}[language=SQL]
SELECT t.id,
       bioseq.score(t.id),
       bioseq.target_start(t.id),
       bioseq.target_end(t.id),
       bioseq.cigar(t.id)
FROM targets AS t
WHERE mmsearch(
        t.sequence,
        $1::bioseq.seq('dna5', 64))
\end{lstlisting}

The projected CIGAR value records the edit path of the current alignment
\cite{sam_format_specification}. The argument of each projection helper
identifies the row lineage to which the value belongs. It does not contain the
score or coordinates. These fields are supplied by the custom executor that
owns the corresponding \texttt{mmsearch} expression.

Let input row $i$ produce $m_i$ accepted matches. The result cardinality is

\[
  N_{\mathrm{out}} = \sum_{i=1}^{N_{\mathrm{in}}} m_i,
  \qquad m_i \geq 0.
\]

An input without a match produces no output. An input with one match behaves
like a conventional filtering result. An input with several matches is
projected several times, once for each match. Table~\ref{tab:scalar-custom-scan}
summarizes the resulting difference in execution semantics.

\begin{table}[htbp]
  \centering
  \footnotesize
  \caption{Scalar evaluation and custom-scan execution.}
  \label{tab:scalar-custom-scan}
  \begin{tabularx}{\linewidth}{
      @{}>{\raggedright\arraybackslash}p{0.22\linewidth}
      >{\raggedright\arraybackslash}X
      >{\raggedright\arraybackslash}X@{}}
    \toprule
    Property & Scalar SQL function & BioSeqDB owner \\
    \midrule
    Work unit & One tuple and one call & A bounded batch of sequence pairs \\
    Cardinality & One scalar value per call & Zero, one, or many rows per input \\
    Alignment fields & Independent calls have no shared match identity &
      All projected fields come from one match record \\
    Execution route & PostgreSQL invokes the function directly &
      The scan selects the CPU or GPU pipeline for the batch \\
    \bottomrule
  \end{tabularx}
\end{table}

The SQL expression therefore acts as a rendezvous between PostgreSQL planning
and BioSeqDB execution. It gives the parser a normal function signature and
gives the planner a recognizable operation. It is not the location at which
alignment is computed.

\section{SQL surface and ownership contract}
\label{sec:mmsearch-sql-rendezvous}

BioSeqDB provides a two-argument \texttt{mmsearch} form that uses session
settings and a three-argument form that accepts typed alignment options. Both
return \texttt{boolean} at the SQL type level. Match projections include the
    dynamic-programming score, query and target coordinates, edit count,
mapping quality, strand, primary status, divergence, seed counts, and CIGAR.
The usual rules for SQL types and null values remain applicable.

The Boolean return type permits \texttt{mmsearch} to appear in a
\texttt{WHERE} or \texttt{JOIN ON} clause, but the ordinary function body is a
failure boundary. If execution reaches it without a BioSeqDB owner, it reports
that the query shape is unsupported. Returning \texttt{true} would be unsafe
because any plan that missed the custom rewrite would accept every row. A
cheap scalar fallback would also lose batching and multiple-match semantics.

Projected helpers follow the same ownership rule. During planning they remain
ordinary expressions, so PostgreSQL can place them below an aggregate or on
the nullable side of a join. During custom execution they are replaced with
typed values associated with one match. A helper attributed to another target
is not filled from the current alignment. This check prevents a score from one
search from being combined with the coordinates of another search in the same
query.

The requested projections also define the work required from the alignment
pipeline. For example, CIGAR output requires traceback data, while a query that
projects only coordinates need not format a CIGAR string. BioSeqDB snapshots
the per-call or session options during planning. At scan initialization, it
combines them with a field mask derived from the planned expressions. The batch
then uses one consistent configuration.

\section{Planner support, custom paths, and custom plans}
\label{sec:planner-support-and-ownership}

The \texttt{mmsearch\_support} function participates in PostgreSQL's planner
support interface \cite{postgresql17_support_functions}. It assigns a high
startup cost to direct scalar evaluation. This cost discourages a plan that
treats alignment as an inexpensive row filter, but it does not by itself
create an owner. Correctness cannot depend on a custom path winning a cost
comparison against an executable scalar path.

Relation and join path hooks perform the structural work. They locate the
exact function expression and its \texttt{RestrictInfo}, then check argument
types, relation lineage, parameterization, option shape, and expression
placement. A base-table predicate can become a \texttt{BioseqScan}. A search
between two relations can become a \texttt{BioseqBatchJoin}. A transformed or
lazy sequence expression may require a \texttt{BioseqSliceProject} below the
owner so that the evaluated sequence has a stable source identity.

Only after these checks does the planner construct a \texttt{CustomPath}. Its
cost estimates source reading, operand evaluation, and alignment. Join
ownership also includes materializing the query side and reading the target
side. These are finite nonnegative costs.
The large scalar cost identifies an undesirable unowned call, while the custom
cost describes the actual work of the owning plan.

Figure~\ref{fig:mmsearch-plan-flow} shows the division of responsibility.

\begin{figure}[htbp]
  \centering
  \footnotesize
  \fbox{\parbox{0.76\linewidth}{\centering
    Typed \texttt{mmsearch} and projection expressions}}
  \\[0.25em]
  $\Downarrow$
  \\[0.25em]
  \fbox{\parbox{0.76\linewidth}{\centering
    Planner hooks verify the call, lineage, options, and placement}}
  \\[0.25em]
  $\Downarrow$
  \\[0.25em]
  \fbox{\parbox{0.76\linewidth}{\centering
    Base or parameterized path uses \texttt{BioseqScan}\\
    Two-child inner search uses \texttt{BioseqBatchJoin}}}
  \\[0.25em]
  $\Downarrow$
  \\[0.25em]
  \fbox{\parbox{0.76\linewidth}{\centering
    \texttt{PlanCustomPath} removes executable copies and records
    serializable owner metadata}}
  \\[0.25em]
  $\Downarrow$
  \\[0.25em]
  \fbox{\parbox{0.76\linewidth}{\centering
    The custom executor batches inputs, injects match fields, and emits rows}}
  \caption{Planning and execution of an owned sequence search. The scalar
  function is retained for typing and planning, not as an alignment route.}
  \label{fig:mmsearch-plan-flow}
\end{figure}

PostgreSQL separates a \texttt{CustomPath} from its executable
\texttt{CustomScan} \cite{postgresql17_custom_scan}. BioSeqDB uses this
boundary to keep process-local state out of the plan. Child paths are carried
as PostgreSQL paths. Expression metadata that must pass through the setrefs
stage is stored as expression nodes. Relation identities, source recipes, and
version tags are stored in copyable plan data. Pipeline objects, buffers, and
sequence providers are created later by the executor.

\texttt{PlanCustomPath} removes the owned \texttt{mmsearch} from ordinary,
vectorized, and storage-level qualification lists. Leaving any copy in those
locations would provide a second execution site for the same operation. The
planner also publishes hidden child outputs before setrefs when an upper node
will need an evaluated slice or a parameter. This allows PostgreSQL to bind
the value through the normal plan tree rather than through hidden global
state.

Plain \texttt{EXPLAIN} initializes the scan only for description and submits
no alignment work.
Prepared statements likewise store expressions and source descriptions, then
bind their values when the plan is executed.

\section{Typed placeholder injection}
\label{sec:match-placeholder-injection}

PostgreSQL evaluates a projection from an expression tree. After planner
setrefs have established the final expression locations, a tree mutator
replaces each attributable helper with a typed PostgreSQL \texttt{Const} node.

The node address is stable throughout execution. Before one output row is
projected, the scan writes the current match fields into the node's value and
null flag. Coordinates, lengths, and edit count use \texttt{BIGINT}.
Scores and counters whose source representation is 32 bit use \texttt{INT4}.
Status fields use \texttt{BOOLEAN}, divergence uses a floating type, and CIGAR
uses \texttt{TEXT}. All placeholders for a row are populated from the same
match record.

The injected node is an executor data channel, not a constant in SQL
semantics. Injection occurs after planning, so PostgreSQL cannot fold the
changing value into the plan. For each emitted alignment, the custom scan
updates the stable nodes, rebuilds the projection state, and then calls
\texttt{ExecProject}.

Joins require an additional placement mechanism. PostgreSQL may move a helper
into a child target list, and an outer join may later null-extend that child.
The projection support function can wrap the helper in a
\texttt{PlaceHolderVar}, preserving its evaluation level, relation set, and
nulling relations \cite{postgresql17_placeholdervar}. The custom scan replaces
the owned expression inside that placement structure. A stable owner identity
distinguishes otherwise identical helper functions that belong to different
searches.

\section{Bounded batch execution}
\label{sec:customscan-batch-state-machine}

The executor alternates between collecting input and draining output. The
batch has a fixed capacity, so the complete relation need not be retained.
When the capacity is reached, or when the source ends, the current batch is
submitted to the selected alignment pipeline. Each input obtains a cursor
over its accepted matches. PostgreSQL then pulls projected rows until these
cursors are empty.

\begin{algorithm}[htbp]
  \caption{Conceptual execution of an owned \texttt{mmsearch}}
  \label{alg:customscan-write-read}
  \footnotesize
  \begin{enumerate}
    \item Read source rows, apply ordinary residual conditions, and capture
      stable sequence references for non-null operands.
    \item When the batch is full or the source ends, run the selected
      alignment pipeline and create a match cursor for each input.
    \item For the next match, restore its input row, populate the requested
      placeholders, and call \textsc{Project}.
    \item Return one tuple, then continue step 3 until all match cursors are
      empty. Release the batch and resume step 1 if input remains.
  \end{enumerate}
\end{algorithm}

This order reconciles batching with PostgreSQL's pull executor. The parent
still asks for one tuple at a time, but expensive processing occurs once for a
group of inputs. An input with no cursor disappears. An input with three cursor
entries is restored and projected three times with three coherent sets of
match fields.

Tuple lifetime is important during this process. A table scan may reuse its
active slot on the next call, and PostgreSQL resets short-lived memory between
returned rows. BioSeqDB therefore copies the row representation needed by the
batch and retains stable sequence-source descriptors. Large or fragmented
sequences can remain locator backed while the provider acquires the regions
required for the current dispatch. If a lazy value must be materialized for
projection, the stable source is restored before another match from the same
input is emitted.

Scalar and vectorized table access converge at the same submission boundary.
A vectorized filter may first produce a keep mask, but every retained row is
then evaluated for alignment operands and submitted to the common batch. This
separation allows ordinary relational filtering to remain vectorized without
turning alignment into a Boolean vector predicate.

An optional prefilter can reject candidates only under its stated lower-bound
condition. Survivors still pass through alignment. A condition on an exact
match field, such as an edit-count limit, is applied to the resulting match
record before projection. The prefilter therefore does not replace the user's
SQL condition with an approximation.

\section{Join execution and rescans}
\label{sec:mmsearch-join-ownership}

A search predicate below a join can remain owned by its base
\texttt{BioseqScan}. This is the simplest case when the query sequence is a
constant or comes from the same row. For correlated nested execution, the
query operand arrives as a PostgreSQL execution parameter. On rescan, the node
clears retained rows and match cursors, rebinds the parameter, and starts a new
batch. Rebinding is required because each outer row may provide a different
query sequence.

A join-level search between target and query relations uses
\texttt{BioseqBatchJoin}. The query child is materialized into stable sequence
references, while the target child is streamed. The owner forms batches of
target and query pairs, performs alignment once, and emits the ordinary columns
of both children together with the current match fields. Residual join
conditions are evaluated at the placement required by PostgreSQL.

When an upper owner needs a lazy slice or another supported transformation, a
projection provider evaluates it. The provider records the base lineage and
the transformation recipe.
The batch join and its projection helpers can then refer to the same evaluated
value without repeating the transformation.

The join-level rewrite is intentionally limited to supported inner-join
shapes with one attributable search operation. The planner declines ownership
when it cannot preserve row security, volatility, nulling, or expression
order. Such a call reaches the explicit unsupported-query boundary if no lower
owner exists. This behavior is preferable to executing a function that cannot
produce the required match rows.

\section{User-visible plan information}
\label{sec:customscan-explain}

The owning node is visible in PostgreSQL \texttt{EXPLAIN}. A base search is
reported as \texttt{Custom Scan (BioseqScan)}, while a two-child search is
reported as \texttt{Custom Scan (BioseqBatchJoin)}. A lazy transformation may
appear as \texttt{BioseqSliceProject}. The owned \texttt{mmsearch} must not
remain as an ordinary \texttt{Filter} or \texttt{Join Filter}, because either
label would indicate an executable scalar copy.

Plain \texttt{EXPLAIN} reports planned information such as projected columns,
operand ownership, source policy, and the planned pipeline architecture.
\texttt{EXPLAIN ANALYZE} adds the executor's observed row counts and the
architecture that actually ran. It can also report batch submissions, index
read or recomputation decisions, storage acquisition, and fallback reasons.
Text, JSON, XML, and YAML output use the same structured description rather
than separate algorithm names maintained by hand.

Ordinary residual conditions can still appear below the custom node, provided
that they do not contain the owned search. Under \texttt{ANALYZE}, the
custom-node row count is the number of emitted matches rather than the number
of source rows. It can therefore exceed the table row count when one target
has several accepted alignments. Conversely, a scanned target with no match
does not contribute an output row. This interpretation follows directly from
the cardinality equation at the start of the chapter.

A useful plan review therefore asks three questions. First, is the expected
BioSeqDB custom node present at the correct relation or join level. Second, is
the semantic \texttt{mmsearch} absent from descendant filters and join
filters. Third, under \texttt{ANALYZE}, do emitted rows and the executed
pipeline agree with the query and configuration. Cost alone does not answer
these questions.

The custom interface has deliberate limits. Batching delays the first result
until at least one dispatch completes. A join owner may need to materialize
one child. Unsupported security or join shapes are rejected. Each new match
projection requires a SQL type, an ownership mapping, and a population rule.
These restrictions follow from the requirement that one search operation has
one executor owner and that every emitted field refers to one match.

\section{Conclusion}
\label{sec:customscan-conclusion}

BioSeqDB keeps a conventional SQL surface while assigning alignment to a
custom relational operator. Planner support marks direct scalar evaluation as
expensive. Structural hooks build an owning custom path, and plan construction
removes every executable copy of the search predicate. Typed placeholders
then connect one match record to ordinary PostgreSQL projection expressions.

The executor collects bounded batches and emits zero, one, or many rows for
each input. Base scans, parameterized rescans, and inner batch joins share this
ownership model. EXPLAIN exposes the selected owner and distinguishes planned
from executed pipeline information. The next chapter describes the physical
sequence representation that allows these nodes to retain small locators and
acquire packed sequence regions only when alignment requires them.
